%--------------------------------------------------------------
\section{v2 改訂点（v1 $\to$ v2、2026年11月22日）}\label{sec:rev}
%--------------------------------------------------------------
v1 の引用の再確認（arXiv ページ・出版社ページ・原文の再読）で見つかった誤りと更新点。数学的な主張（定理・命題）は変更していない。
\begin{enumerate}
\item \textbf{arXiv:0905.3227 の著者}：Gualtieri ではなく \textbf{Hitchin}（\emph{Poisson modules and generalized geometry}）。v1 の参考文献と \S\ref{sec:ledger} の記述を訂正。
\item \textbf{Gualtieri 2011 の番号}：楕円性は Prop.\ 3.13 と Cor.\ 3.14（v1 の「Prop.\ 3.16」は存在しない）。Def.\ 3.18 は Picard 群の定義で、一般化正則束は Thm.\ 3.15 の後に散文で導入される。一般化 Darboux は Thm.\ 4.3（v1 の「4.35」は誤り）。v1 の「例 4.1--4.2」は存在せず、Def.\ 4.1 が正則点の定義。
\item \textbf{Pal--Poddar（arXiv:2404.18113）}：題名 \emph{Strong generalized holomorphic principal bundles}、著者 D.\ Pal と M.\ Poddar。v1 の「著者表示の食い違い」は再現しなかった。定理番号を整理（Hodge 分解・Serre 双対性は Thm.\ 11.2, 11.4, 11.6, 11.7、消滅は Thm.\ 11.12）。Thm.\ 11.12 の主張と、「消滅は横断次数についてのみ」という本ノートの読みとの整合を、\S\S5--11, 13 の再読で確認（群は同一ではない）。原文の疑義の原因（特性多項式の定義の食い違い）を更新。
\item \textbf{Lang--Jia--Liu}：arXiv:2102.07177、\emph{The Atiyah class of generalized holomorphic vector bundles}。局所自明化で定義される GH 束の論文で、Pal--Poddar の SGH 束とは別。
\item \textbf{Hu（Tengzhou Hu）2403.10876}：印刷された消滅定理の証明に符号のずれがあることを本文の再読で確認した（主張自体は古典と整合）。
\item \textbf{書誌の確定}：Hitchin (Moscow Math.\ J.\ 12 (2012) 567--591)、Pym (Lett.\ Math.\ Phys.\ 108 (2018) 573--632)、Hausel--Hunsicker--Mazzeo (Duke 122 (2004) 485--548)、Timmerscheidt (J.\ reine angew.\ Math.\ 379 (1987) 152--171、v1 の記憶は誤り)、Mokri の年・DOI、Goto の arXiv 番号 1501.03398、Arapura の arXiv 番号。
\item \textbf{Bailey--Gualtieri Thm.\ 7.1}：実 Poisson 構造の積分可能性が必要十分であることを確認し、Hausdorff・$s$-連結性に関する注意を追記。
\end{enumerate}
なお、本節で確認済みとした以外の \tQ\ 項目は、\S\ref{sec:ledger} に残している。
